\section{Введение}

\subsection{Наименование разработки}

\subsubsection{Наименование разработки на русском языке}

<<Генератор данных для аналитического бенчмарка СУБД TPC-H>>.

\subsubsection{Наименование разработки на английском языке}

<<Data Generator for the TPC-H Analytical DBMS Benchmark>>.

\subsubsection{Условное обозначение темы разработки}

\texttt{Schengen}.

Исходный код программы размещается в репозитории \texttt{Schengen}
\srcref{10}.

\subsection{Документы, являющиеся основанием для разработки}

Разработка ведется на основании следующих документов:
\begin{enumerate}
    \item учебный план подготовки бакалавров по направлению 09.03.04
    <<Программная инженерия>> федерального государственного автономного
    образовательного учреждения высшего образования Национальный
    исследовательский университет <<Высшая школа экономики>>;
    \item утвержденная академическим руководителем образовательной программы
    тема выпускной квалификационной работы <<Генератор данных для аналитического
    бенчмарка СУБД TPC-H>>;
    \item техническое задание \texttt{RU.17701729.04.03-01 ТЗ 01-1},
    разработанное в рамках выполнения выпускной квалификационной работы.
\end{enumerate}

\subsection{Предпосылки разработки}

Необходимость разработки определяется тем, что существующие генераторы данных
TPC-H \srcref{1} либо ориентированы на вывод в строковые файлы
\texttt{tbl}/\texttt{csv} \srcref{4}, либо не учитывают требований современной
инфраструктуры аналитического хранения и обработки данных, использующей
колонoчные форматы \texttt{Parquet} \srcref{5} и \texttt{ORC} \srcref{6}, а
также табличные слои lakehouse-платформ \srcref{7}.

В практических сценариях разработки и исследования аналитических систем
требуется средство, обеспечивающее детерминированную генерацию таблиц TPC-H
\srcref{1}, прямую запись в современные колонoчные форматы \texttt{Parquet}
\srcref{5}, \texttt{ORC} \srcref{6}, \texttt{Lance} \srcref{8} и
\texttt{Vortex} \srcref{9}, поддержку параллельного формирования партиций в
духе систем параллельной обработки данных \srcref{2} и возможность работы как
с локальной файловой системой, так и с \texttt{S3}-совместимыми хранилищами.
